\chapter{Loop anatomy}
\label{ch:loop}

\section*{What this chapter is for}

A hiring process is a sequence of filters and they do not measure the same
thing. Knowing which filter you are in changes what a good answer looks like,
and the most common preparation error is to prepare for all of them the same
way \dash{} usually by preparing for the technical ones, which are not where
most candidates are lost.

\section{The shape}

Disclosure is rare, so the numbers below come from the one aggregation with a
stated method: \reported{a median of four stages, with most companies running
three to five, and a range from two to seven; end to end, three to six rounds
over two to six weeks.}{field-guide-process} The same corpus reports
\reported{a take-home in about a third of the loops that disclose one, typically
two to three hours, and usually replacing an on-site round rather than adding
to the load.}{field-guide-process}

A typical order:

\begin{kittable}{@{}lXl@{}}
\textbf{Stage} & \textbf{What it actually measures} & \textbf{Length} \\
\midrule
Recruiter screen & can you be hired: band, location, notice, coherence & 15--30 min \\
Technical & one of: live coding, system design, code review & 45--60 min \\
Hiring manager & will this person work here, on this & 45--60 min \\
Take-home or live build & how you work when nobody is prompting you & 2--3 h \\
Project deep-dive & what your scope actually was & 45--60 min \\
Behavioural / values & how you behave when it goes wrong & 45--60 min \\
Executive & a final judgement, on whatever they choose & 15--60 min \\
\end{kittable}

\section{What ends you at each}

This is the useful column and it is rarely the technical one.

\textbf{The screen.} Almost nothing technical. It ends on a band mismatch, a
location or work-model mismatch, a notice period, or an answer that does not
match the written application. \judgement{It is the least demanding stage and
frequently the one with the lowest pass rate, which is a strange combination
worth taking seriously: the preparation that pays here is thirty minutes of
consistency work, not a week of study.}

\textbf{The technical round.} Ends on not being able to explain your own
reasoning, far more often than on not knowing something. The same corpus notes
that interviewers watch for reasoning rather than
output.\source{field-guide}

\textbf{The hiring manager.} Ends on a mismatch of what the job is. This is the
stage where the segment question from Chapter~\ref{ch:segment} gets settled,
and where a candidate who prepared for the wrong one of the three jobs finds
out.

\textbf{The build round.} Ends on scope: not finishing anything, rather than
not finishing everything. Part~\ref{part:build} is about that.

\textbf{The project deep-dive.} Ends on the scope of the work you describe.
Chapter~\ref{ch:deep-dive} is the whole of it, because it is where a Staff
decision is actually made.

\textbf{The executive round.} Ends on whatever that person is worried about.
This is the stage where an objection that three people accepted can be raised
again by somebody who is not bound by their acceptance, and
Chapter~\ref{ch:calibration} is about why your estimate has to allow for it.

\section{The two formats you must ask about}

One question in the screen changes your preparation more than any other:
\emph{is AI permitted in the technical rounds?}

The answer is genuinely split. \reported{One industry survey reports that
roughly half of companies permit it, broadly or with constraints, a fifth decide
case by case, and the remainder do not.}{coderpad-hiring-2026} At the same time
\reported{several employers ban it explicitly, and in-person interviewing has
been rising rather than falling.}{field-guide-trends}

So there is no default to assume, and Chapter~\ref{ch:formats} is about
preparing for both without doubling the work.

\section{What to ask, and when}

\begin{kitmethod}
\textbf{In the screen:} how many stages, what each one is, whether there is a
take-home or a live build, whether AI is permitted in it, and who you will
speak to at each. All of it is administrative, all of it is reasonable to ask,
and the answers reshape your week.

\textbf{Before a technical round:} what format, and roughly what surface.
``Should I expect this to be closer to system design or closer to a working
session?'' is a fair question and the answer is usually given.

\textbf{Before the build round:} the length, the tooling, whether you choose
the stack, and whether they expect something running at the end. If a brief
exists, ask for it. Part~\ref{part:build} assumes you have these answers.
\end{kitmethod}

\begin{kitnote}
Some of these questions go unanswered, and that is information too. A process
that cannot say what its own stages are is a process being assembled as it
runs, and it will be assembled around you.
\end{kitnote}

\begin{drill}
Write down the loop for the process you are in, as a table with one row per
stage, filling in what you know and marking the gaps. Then send the gaps to
whoever is coordinating, in one message. Most of them will come back answered,
and the ones that do not are the ones to plan around.
\end{drill}

\section*{What this chapter does not claim}

The stage counts describe the small subset of companies that publish their
process \dash{} around fifty, in the corpus behind them. That is a real
measurement of a self-selected group and not a description of the market.

The ``what ends you'' column is the author's judgement, assembled from what
each stage is structurally able to observe. It is a useful prior and it is not
a measurement, and no figure is offered for pass rates at any stage: those
circulate widely and are refused in Appendix~\ref{app:refused}.
